\documentclass[12pt,a4paper]{article}

% --- Packages ---
\usepackage[utf8]{inputenc}
\usepackage[T1]{fontenc}
\usepackage[english]{babel}
\usepackage{geometry}
\usepackage{graphicx}
\usepackage{hyperref}
\usepackage{listings}
\usepackage{xcolor}
\usepackage{booktabs}
\usepackage{longtable}
\usepackage{enumitem}
\usepackage{fancyhdr}
\usepackage{titlesec}
\usepackage{caption}
\usepackage{float}
\usepackage{parskip}
\usepackage{tocloft}

% --- Page Geometry ---
\geometry{margin=2.5cm}

% --- Hyperlinks ---
\hypersetup{
    colorlinks=true,
    linkcolor=blue!60!black,
    urlcolor=blue!60!black,
    citecolor=blue!60!black,
    pdftitle={DMS -- Diabetes Management System Documentation},
    pdfauthor={DMS Development Team},
}

% --- Code Listing Style ---
\definecolor{codebg}{HTML}{F5F5F5}
\definecolor{codeframe}{HTML}{CCCCCC}
\definecolor{codegreen}{HTML}{34A853}
\definecolor{codepurple}{HTML}{9334E6}
\definecolor{codeblue}{HTML}{1A73E8}

\lstdefinestyle{dartstyle}{
    backgroundcolor=\color{codebg},
    frame=single,
    rulecolor=\color{codeframe},
    basicstyle=\ttfamily\footnotesize,
    keywordstyle=\color{codeblue}\bfseries,
    stringstyle=\color{codegreen},
    commentstyle=\color{gray}\itshape,
    breaklines=true,
    showstringspaces=false,
    tabsize=2,
    numbers=left,
    numberstyle=\tiny\color{gray},
    numbersep=5pt,
}

\lstdefinestyle{yamlstyle}{
    backgroundcolor=\color{codebg},
    frame=single,
    rulecolor=\color{codeframe},
    basicstyle=\ttfamily\footnotesize,
    keywordstyle=\color{codepurple},
    stringstyle=\color{codegreen},
    breaklines=true,
    showstringspaces=false,
    tabsize=2,
}

\lstdefinestyle{pythonstyle}{
    backgroundcolor=\color{codebg},
    frame=single,
    rulecolor=\color{codeframe},
    basicstyle=\ttfamily\footnotesize,
    keywordstyle=\color{codeblue}\bfseries,
    stringstyle=\color{codegreen},
    commentstyle=\color{gray}\itshape,
    breaklines=true,
    showstringspaces=false,
    tabsize=4,
    numbers=left,
    numberstyle=\tiny\color{gray},
    numbersep=5pt,
}

% --- Headers & Footers ---
\pagestyle{fancy}
\fancyhf{}
\fancyhead[L]{\small DMS -- Diabetes Management System}
\fancyhead[R]{\small v1.0.0}
\fancyfoot[C]{\thepage}
\renewcommand{\headrulewidth}{0.4pt}

% --- Title ---
\title{
    \vspace{-1cm}
    {\Huge\bfseries DMS}\\[0.3cm]
    {\LARGE Diabetes Management System}\\[0.5cm]
    {\large Technical Documentation}\\[0.3cm]
    \rule{\textwidth}{0.4pt}
}
\author{DMS Development Team}
\date{February 2026 \\ Version 1.0.0}

% ==============================================================================
\begin{document}
% ==============================================================================

\maketitle
\thispagestyle{empty}

\vfill
\begin{center}
\textit{A comprehensive mobile application for continuous glucose monitoring,\\
diabetes event tracking, and patient--doctor collaboration.}
\end{center}
\vfill

\clearpage
\tableofcontents
\clearpage

%==============================================================================
\section{Introduction}
%==============================================================================

\subsection{Project Overview}

DMS (Diabetes Management System) is a cross-platform mobile application built with Flutter that enables real-time glucose monitoring through continuous glucose monitoring (CGM) sensors, specifically the Dexcom G7 platform. The application supports two distinct user roles---\textbf{Patients} and \textbf{Doctors}---and provides a comprehensive ecosystem for diabetes data management, analysis, and collaboration.

The system consists of three main components:
\begin{enumerate}
    \item \textbf{Flutter Mobile Application} -- The primary user-facing application supporting Android, iOS, Windows, Linux, macOS, and web platforms.
    \item \textbf{FastAPI Backend Server} -- A Python-based REST API providing Dexcom OAuth2 integration, data synchronisation, and AI-powered glucose analysis.
    \item \textbf{Firebase Cloud Services} -- Authentication, real-time database (Cloud Firestore), and push notification infrastructure.
\end{enumerate}

\subsection{Purpose and Motivation}

Continuous glucose monitoring generates vast amounts of data that can be overwhelming for patients to interpret on their own. DMS addresses this by:
\begin{itemize}
    \item Providing clear, interactive visualisations of glucose trends over configurable time ranges (3, 6, 12, and 24 hours).
    \item Offering AI-powered analysis of glucose patterns with personalised recommendations.
    \item Enabling secure data sharing between patients and their healthcare providers.
    \item Tracking diabetes-related events (insulin doses, meals, physical activity) alongside glucose readings.
    \item Delivering timely alerts for hypo- and hyperglycaemic events.
\end{itemize}

\subsection{Medical Disclaimer}

\begin{quote}
\textbf{IMPORTANT:} DMS uses an unofficial Dexcom Share API and should be used for \textit{informational purposes only}. Do \textbf{not} use this application for critical medical treatment decisions. Always consult your healthcare provider before making changes to your diabetes management plan.
\end{quote}

%==============================================================================
\section{System Architecture}
%==============================================================================

\subsection{High-Level Architecture}

The DMS platform follows a three-tier architecture:

\begin{verbatim}
+-------------------+       +-------------------+       +-------------------+
|   Flutter App     | <---> |   FastAPI Backend  | <---> |   Dexcom Share    |
|  (Mobile/Desktop) |       |   (Python 3.12+)  |       |   API (HTTPS)     |
+-------------------+       +-------------------+       +-------------------+
        |                           |
        v                           v
+-------------------+       +-------------------+
|  Firebase Auth    |       |   PostgreSQL DB   |
|  Cloud Firestore  |       |  (Token Storage)  |
+-------------------+       +-------------------+
\end{verbatim}

\subsection{Data Flow}

\begin{enumerate}
    \item \textbf{Authentication:} Users authenticate via Firebase Auth (email/password). Credentials are persisted locally using \texttt{SharedPreferences}.
    \item \textbf{Glucose Data (Real Sensor):} The Flutter app authenticates directly with the Dexcom Share API using the patient's Dexcom credentials, then polls for glucose readings at regular intervals.
    \item \textbf{Glucose Data (Mock/Demo):} For demonstration purposes, mock glucose data is generated algorithmically and stored in Cloud Firestore.
    \item \textbf{Doctor Access:} Doctors query Firestore for their assigned patients' glucose readings and diabetes events. Firestore security rules enforce the doctor--patient relationship.
    \item \textbf{AI Analysis:} The Flutter app sends glucose data to the FastAPI backend, which forwards it to a large language model (LLM) for analysis and returns personalised recommendations.
\end{enumerate}

\subsection{Technology Stack}

\begin{longtable}{p{4cm} p{3.5cm} p{7cm}}
\toprule
\textbf{Layer} & \textbf{Technology} & \textbf{Purpose} \\
\midrule
\endhead
Frontend Framework & Flutter 3.9+ & Cross-platform UI (Android, iOS, Windows, Web) \\
Programming Language & Dart 3.9+ & Frontend application logic \\
State Management & Provider 6.1 & Reactive state management with \texttt{ChangeNotifier} \\
Navigation & GoRouter 17.0 & Declarative routing with deep linking support \\
Charts & fl\_chart 1.1 & Interactive glucose trend line charts \\
Typography & Google Fonts & Roboto font family (Material Design 3) \\
Authentication & Firebase Auth 6.1 & Email/password authentication \\
Database & Cloud Firestore 6.1 & NoSQL document database for user/health data \\
CGM Integration & Dexcom Share API & Real-time glucose readings from Dexcom G7 sensors \\
HTTP Client & http 1.2 & Direct REST API calls (Dexcom bypass, AI backend) \\
Local Storage & SharedPreferences & Persisting credentials, settings, session state \\
Notifications & flutter\_local\_notifications & Local push notifications for glucose alerts \\
Markdown & flutter\_markdown & Rendering AI assistant responses \\
Backend Framework & FastAPI 0.115 & Python REST API server \\
Backend ORM & SQLAlchemy 2.0 & Database abstraction for PostgreSQL \\
Backend Auth & Firebase Admin SDK & Server-side token verification \\
Encryption & cryptography 44.0 & AES encryption for stored OAuth tokens \\
Scheduling & APScheduler 3.10 & Background job scheduling for data sync \\
\bottomrule
\end{longtable}

%==============================================================================
\section{Project Structure}
%==============================================================================

\subsection{Flutter Application}

\begin{lstlisting}[style=yamlstyle, title={lib/ directory structure}]
lib/
  main.dart                        # App entry point, Firebase init
  firebase_options.dart            # Firebase configuration (auto-generated)
  core/
    constants/
      app_constants.dart           # Glucose thresholds, time ranges, keys
    router/
      app_router.dart              # GoRouter route definitions
    theme/
      app_theme.dart               # Material 3 light/dark themes
    config.dart                    # AI backend URL configuration
  features/
    auth/screens/
      splash_screen.dart           # Initial loading / auth check
      login_screen.dart            # Email/password login
      register_screen.dart         # New account registration
      role_selection_screen.dart   # Patient or Doctor role picker
    patient/screens/
      patient_dashboard_screen.dart # Main patient view (glucose + chart)
      patient_history_screen.dart   # Historical glucose data
      patient_connections_screen.dart # Doctor connections
      patient_profile_screen.dart   # Patient profile settings
      add_event_screen.dart         # Log insulin, meals, activity
      ai_assistant_screen.dart      # AI-powered glucose analysis
    doctor/screens/
      doctor_dashboard_screen.dart  # Patient list + search + alerts
      doctor_patients_screen.dart   # Full patient directory
      patient_detail_screen.dart    # Individual patient data view
    settings/screens/
      settings_screen.dart          # App settings (theme, units, etc.)
      alerts_settings_screen.dart   # Glucose alert thresholds
      dexcom_connection_screen.dart # Dexcom account linking
      notification_settings_screen.dart # Push notification config
  models/
    user.dart                      # User model (Patient/Doctor)
    glucose_reading.dart           # Glucose reading from sensor
    diabetes_event.dart            # Events: insulin, meal, activity, note
  providers/
    auth_provider.dart             # Authentication state
    glucose_provider.dart          # Glucose data management
    events_provider.dart           # Diabetes events management
    settings_provider.dart         # App preferences
    ai_assistant_provider.dart     # AI chat state management
  services/
    firebase_auth_service.dart     # Firebase Auth wrapper
    firestore_service.dart         # Firestore CRUD operations
    dexcom_service.dart            # Dexcom Share API integration
    ai_assistant_service.dart      # LLM-backed analysis service
    notification_service.dart      # Local notification scheduling
    mock_data_service.dart         # Demo data generation
    mock_analysis_service.dart     # Fallback AI analysis
    data_sync_service.dart         # Cloud data synchronisation
  widgets/
    glucose_chart.dart             # fl_chart glucose line chart
\end{lstlisting}

\subsection{Backend Server}

\begin{lstlisting}[style=yamlstyle, title={backend/ directory structure}]
backend/
  main.py                          # Uvicorn entry point
  Dockerfile                       # Docker containerisation
  docker-compose.yml               # Multi-service orchestration
  requirements.txt                 # Python dependencies
  alembic.ini                      # Database migration config
  alembic/
    env.py                         # Alembic environment
    script.py.mako                 # Migration template
  app/
    __init__.py
    main.py                        # FastAPI app creation + lifespan
    config.py                      # Pydantic settings (env vars)
    database.py                    # SQLAlchemy engine + session
    models/
      database.py                  # ORM models (tokens, sessions)
      dexcom.py                    # Dexcom API response schemas
    routers/
      health.py                    # Health check endpoints
      dexcom.py                    # Dexcom OAuth + data endpoints
      sync.py                      # Data synchronisation endpoints
    services/
      dexcom_api.py                # Dexcom REST API client
      dexcom_scraper.py            # Alternative data fetcher
      firestore_service.py         # Server-side Firestore client
      job_scheduler.py             # APScheduler background jobs
      rate_limiter.py              # Dexcom API rate limiting
      token_storage.py             # Encrypted OAuth token storage
    utils/                         # Utility functions
\end{lstlisting}

%==============================================================================
\section{Data Models}
%==============================================================================

\subsection{User Model}

The \texttt{User} model represents both patients and doctors, distinguished by the \texttt{role} field.

\begin{longtable}{p{3.5cm} p{3cm} p{8cm}}
\toprule
\textbf{Field} & \textbf{Type} & \textbf{Description} \\
\midrule
\endhead
\texttt{id} & \texttt{String} & Firebase Auth UID (primary key) \\
\texttt{email} & \texttt{String} & User email address \\
\texttt{name} & \texttt{String} & Display name \\
\texttt{role} & \texttt{String} & \texttt{'patient'} or \texttt{'doctor'} \\
\texttt{profileImageUrl} & \texttt{String?} & Avatar URL (optional) \\
\texttt{createdAt} & \texttt{DateTime?} & Account creation timestamp \\
\texttt{specialization} & \texttt{String?} & Doctor-only: medical specialty \\
\texttt{licenseNumber} & \texttt{String?} & Doctor-only: medical license \\
\texttt{doctorId} & \texttt{String?} & Patient-only: assigned doctor's UID \\
\texttt{dateOfBirth} & \texttt{DateTime?} & Patient-only: date of birth \\
\texttt{diabetesType} & \texttt{String?} & Patient-only: e.g.\ \texttt{'Type 1'}, \texttt{'Type 2'} \\
\bottomrule
\end{longtable}

The model provides helper getters: \texttt{isDoctor} and \texttt{isPatient}, as well as \texttt{fromJson}, \texttt{toJson}, and \texttt{copyWith} methods for serialisation and immutable updates.

\subsection{Glucose Reading Model}

Each glucose reading represents a single data point from the CGM sensor or a manual entry.

\begin{longtable}{p{3.5cm} p{3cm} p{8cm}}
\toprule
\textbf{Field} & \textbf{Type} & \textbf{Description} \\
\midrule
\endhead
\texttt{id} & \texttt{String} & Unique reading identifier \\
\texttt{patientId} & \texttt{String} & Patient's Firebase UID \\
\texttt{value} & \texttt{double} & Glucose level in mg/dL \\
\texttt{timestamp} & \texttt{DateTime} & When the reading was taken \\
\texttt{trend} & \texttt{String?} & Trend direction: \texttt{'rising'}, \texttt{'falling'}, \texttt{'stable'}, \texttt{'rising\_fast'}, \texttt{'falling\_fast'} \\
\texttt{source} & \texttt{String} & Data source: \texttt{'sensor'}, \texttt{'manual'}, \texttt{'calibration'} \\
\bottomrule
\end{longtable}

The \texttt{GlucoseStatus} enum classifies readings based on clinical thresholds:
\begin{itemize}
    \item \textbf{Low} ($< 70$ mg/dL) -- Hypoglycaemia risk
    \item \textbf{Normal} ($70$--$180$ mg/dL) -- Target range
    \item \textbf{High} ($180$--$250$ mg/dL) -- Hyperglycaemia
    \item \textbf{Very High} ($> 250$ mg/dL) -- Severe hyperglycaemia
\end{itemize}

\subsection{Diabetes Event Model}

Events track patient actions and observations alongside glucose data.

\begin{longtable}{p{3.5cm} p{3cm} p{8cm}}
\toprule
\textbf{Field} & \textbf{Type} & \textbf{Description} \\
\midrule
\endhead
\texttt{id} & \texttt{String} & Unique event identifier \\
\texttt{patientId} & \texttt{String} & Patient's Firebase UID \\
\texttt{type} & \texttt{EventType} & Category of event (see below) \\
\texttt{timestamp} & \texttt{DateTime} & When the event occurred \\
\texttt{data} & \texttt{Map<String, dynamic>} & Type-specific payload (units, carbs, duration, etc.) \\
\texttt{notes} & \texttt{String?} & Free-text note (optional) \\
\bottomrule
\end{longtable}

Supported \texttt{EventType} values:
\begin{itemize}
    \item \texttt{bloodGlucose} -- Manual finger-prick measurement
    \item \texttt{insulin} -- Insulin dose (units in \texttt{data['units']})
    \item \texttt{meal} -- Carbohydrate intake (grams in \texttt{data['carbs']})
    \item \texttt{activity} -- Physical activity (minutes in \texttt{data['duration']})
    \item \texttt{fastingGlucose} -- Morning fasting measurement
    \item \texttt{note} -- General observation or annotation
\end{itemize}

%==============================================================================
\section{State Management}
%==============================================================================

The application uses the \textbf{Provider} package with \texttt{ChangeNotifier} classes for reactive state management. All providers are registered at the root of the widget tree via \texttt{MultiProvider} in \texttt{main.dart}.

\subsection{AuthProvider}

Manages the full authentication lifecycle:
\begin{itemize}
    \item Listens to Firebase Auth state changes via \texttt{authStateChanges} stream.
    \item On login, fetches the \texttt{User} document from Firestore and caches it.
    \item For the mock account (\texttt{mocked@test.pl}), generates demo data automatically using \texttt{MockDataService}.
    \item Persists session state (\texttt{isLoggedIn}, \texttt{userRole}, \texttt{userId}, \texttt{userName}) to \texttt{SharedPreferences}.
    \item Provides \texttt{login()}, \texttt{register()}, \texttt{logout()}, \texttt{sendPasswordResetEmail()}, and \texttt{checkAuthState()} methods.
\end{itemize}

\subsection{GlucoseProvider}

Manages glucose reading data from multiple sources:
\begin{itemize}
    \item \texttt{initializeDexcom(patientId)} -- Authenticates with Dexcom and loads real sensor data.
    \item \texttt{loadGlucoseReadings(patientId, hours)} -- Fetches readings via the Dexcom service.
    \item \texttt{loadGlucoseReadingsFromFirestore(patientId, hours)} -- Loads data from Cloud Firestore (for mock accounts and doctor views).
    \item \texttt{initializeMockData(patientId)} -- Generates deterministic mock data using a seeded random number generator.
    \item \texttt{getReadingsForTimeRange(hours)} -- Filters cached readings by time window.
    \item \texttt{getStatistics(hours)} -- Computes average, min, max, in-range percentage for a time period.
\end{itemize}

\subsection{EventsProvider}

Handles diabetes event CRUD operations:
\begin{itemize}
    \item \texttt{addEvent(event)} -- Creates a new event and persists to Firestore.
    \item \texttt{loadEvents(patientId)} -- Fetches events from the database.
    \item Maintains a local cache for offline resilience.
\end{itemize}

\subsection{SettingsProvider}

Manages user preferences stored in \texttt{SharedPreferences}:
\begin{itemize}
    \item Dark mode toggle
    \item Notification preferences
    \item Glucose unit selection (mg/dL vs.\ mmol/L)
\end{itemize}

\subsection{AiAssistantProvider}

Manages the AI chat interface state:
\begin{itemize}
    \item Maintains conversation history (list of messages).
    \item Sends glucose data to the AI backend for analysis.
    \item Handles loading states and error recovery.
\end{itemize}

%==============================================================================
\section{Services Layer}
%==============================================================================

\subsection{Firebase Authentication Service}

\texttt{FirebaseAuthService} wraps the Firebase Auth SDK with user-friendly error handling.

\begin{itemize}
    \item \texttt{signInWithEmailAndPassword()} -- Returns a \texttt{User} model on success.
    \item \texttt{registerWithEmailAndPassword()} -- Creates both a Firebase Auth account and a Firestore user document.
    \item \texttt{signOut()}, \texttt{sendPasswordResetEmail()}, \texttt{sendEmailVerification()}, \texttt{deleteAccount()}.
    \item Maps Firebase error codes to human-readable messages (e.g.\ \texttt{'network-request-failed'} $\rightarrow$ \texttt{'Network connection error.'}).
\end{itemize}

\subsection{Firestore Service}

\texttt{FirestoreService} provides all Cloud Firestore CRUD operations:

\begin{longtable}{p{5.5cm} p{9cm}}
\toprule
\textbf{Method} & \textbf{Description} \\
\midrule
\endhead
\texttt{createUser(user)} & Create new user document \\
\texttt{getUserById(userId)} & Fetch user by Firebase UID \\
\texttt{getUserByEmail(email)} & Look up user by email address \\
\texttt{updateUser(user)} & Update existing user document \\
\texttt{getPatientsByDoctorId(id)} & Query all patients assigned to a doctor \\
\texttt{assignPatientToDoctor(pid, did)} & Set the \texttt{doctorId} field on a patient \\
\texttt{getGlucoseReadings(pid, hours)} & Fetch glucose readings for a patient \\
\texttt{getDiabetesEvents(pid, hours)} & Fetch diabetes events for a patient \\
\texttt{getGlucoseReadingsStream(pid)} & Real-time stream of glucose data \\
\texttt{getDiabetesEventsStream(pid)} & Real-time stream of events \\
\texttt{getPatientsStream(doctorId)} & Real-time stream of a doctor's patients \\
\bottomrule
\end{longtable}

\subsection{Dexcom CGM Service}

\texttt{DexcomService} provides integration with Dexcom continuous glucose monitors via the \textbf{Dexcom Share API}. Due to a timestamp-parsing bug in the \texttt{dexcom} Dart package (v1.0.7) that causes \texttt{FormatException} on European timezone offsets (\texttt{+0100}), the service \textbf{bypasses the package} for data fetching and makes direct HTTP calls.

\subsubsection{Authentication Flow}
\begin{enumerate}
    \item \texttt{authenticate(username, password, region)} is called with the user's Dexcom Share credentials.
    \item \texttt{\_createSession()} sends two sequential HTTP POST requests:
    \begin{itemize}
        \item \texttt{/General/AuthenticatePublisherAccount} -- Returns an \texttt{accountId}.
        \item \texttt{/General/LoginPublisherAccountById} -- Returns a \texttt{sessionId}.
    \end{itemize}
    \item The \texttt{sessionId} is cached and used for all subsequent data requests.
    \item If a request returns HTTP~500, the session is automatically refreshed.
\end{enumerate}

\subsubsection{Supported Regions}
\begin{longtable}{p{2cm} p{6cm} p{5cm}}
\toprule
\textbf{Region} & \textbf{Base URL} & \textbf{Application ID} \\
\midrule
\texttt{us} & \texttt{https://share2.dexcom.com/\ldots} & \texttt{d89443d2-327c-\ldots} \\
\texttt{ous} & \texttt{https://shareous1.dexcom.com/\ldots} & \texttt{d89443d2-327c-\ldots} \\
\texttt{jp} & \texttt{https://share.dexcom.jp/\ldots} & \texttt{d89443d2-327c-\ldots} \\
\bottomrule
\end{longtable}

\subsubsection{Timestamp Parsing}

The Dexcom Share API returns timestamps in multiple formats:
\begin{itemize}
    \item Raw integer milliseconds: \texttt{1770654779246}
    \item Milliseconds with timezone offset: \texttt{1770654779246+0100}
    \item .NET \texttt{/Date()} wrapper: \texttt{/Date(1770654779246+0100)/}
\end{itemize}

The custom \texttt{\_parseDexcomTimestamp()} method handles all three formats by stripping the timezone offset with the regex \texttt{\^{}(\textbackslash d+)[+-]\textbackslash d\{4\}\$} and parsing the remaining milliseconds.

\subsubsection{Key Methods}
\begin{itemize}
    \item \texttt{getReadings(patientId, minutes, maxCount)} -- Fetch glucose readings for a time range.
    \item \texttt{getCurrentReading(patientId)} -- Get the most recent glucose value.
    \item \texttt{getFullHistory(patientId)} -- Retrieve up to 7 days of data.
    \item \texttt{streamReadings(patientId, interval)} -- Polling-based stream that yields new readings at a configurable interval (default: 5~minutes).
\end{itemize}

\subsection{AI Assistant Service}

\texttt{AiAssistantService} connects to a custom FastAPI backend that proxies requests to a large language model for glucose analysis.

\begin{itemize}
    \item \texttt{checkHealth()} -- Verifies the backend is online.
    \item \texttt{analyzeGlucose(currentReading, recentReadings, userContext)} -- Sends glucose data and receives personalised advice.
    \item Builds structured prompts including current value, trend, recent history, and optional patient context.
    \item All responses include a medical disclaimer.
    \item Falls back to \texttt{MockAnalysisService} when the backend is unavailable.
\end{itemize}

\subsection{Notification Service}

\texttt{NotificationService} uses \texttt{flutter\_local\_notifications} to deliver alerts:
\begin{itemize}
    \item Glucose high/low threshold alerts.
    \item Sensor disconnection warnings.
    \item Configurable through the notification settings screen.
\end{itemize}

\subsection{Mock Data Service}

\texttt{MockDataService} generates realistic demo data for the mock account (\texttt{mocked@test.pl}):
\begin{itemize}
    \item Generates 7~days of glucose readings at 5-minute intervals.
    \item Simulates meal spikes at 08:00, 13:00, and 19:00.
    \item Values are clamped between 60 and 300~mg/dL with a fixed random seed (\texttt{Random(42)}) for consistency.
    \item Stores generated data in Cloud Firestore.
\end{itemize}

%==============================================================================
\section{Navigation and Routing}
%==============================================================================

The application uses \textbf{GoRouter} for declarative, URL-based navigation. All routes are defined in \texttt{AppRouter}:

\begin{longtable}{p{5cm} p{4.5cm} p{5cm}}
\toprule
\textbf{Path} & \textbf{Screen} & \textbf{Description} \\
\midrule
\endhead
\texttt{/} & SplashScreen & Initial auth check \\
\texttt{/login} & LoginScreen & Email/password login \\
\texttt{/register} & RegisterScreen & Account creation \\
\texttt{/role-selection} & RoleSelectionScreen & Patient/Doctor choice \\
\midrule
\texttt{/patient} & PatientDashboardScreen & Main patient view \\
\texttt{/patient/history} & PatientHistoryScreen & Historical data \\
\texttt{/patient/connections} & PatientConnectionsScreen & Doctor connections \\
\texttt{/patient/profile} & PatientProfileScreen & Profile settings \\
\texttt{/patient/add-event} & AddEventScreen & Log events \\
\texttt{/patient/ai-assistant} & AiAssistantScreen & AI analysis chat \\
\midrule
\texttt{/doctor} & DoctorDashboardScreen & Patient list + alerts \\
\texttt{/doctor/patients} & DoctorPatientsScreen & Full patient directory \\
\texttt{/doctor/patient/:id} & PatientDetailScreen & Patient detail view \\
\midrule
\texttt{/settings} & SettingsScreen & App settings \\
\texttt{/settings/alerts} & AlertsSettingsScreen & Alert thresholds \\
\texttt{/settings/dexcom} & DexcomConnectionScreen & Sensor linking \\
\texttt{/settings/notifications} & NotificationSettingsScreen & Push notifications \\
\bottomrule
\end{longtable}

%==============================================================================
\section{Theming and UI Design}
%==============================================================================

\subsection{Design Language}

The application follows \textbf{Material Design 3} principles, inspired by the Dexcom G7 app design:
\begin{itemize}
    \item Green accent colour (\texttt{\#34A853}) for health/wellness context.
    \item Blue primary colour (\texttt{\#1A73E8}) matching Google's Material palette.
    \item Card-based layouts with 20px border radius.
    \item Bottom navigation bar for primary screen switching.
    \item Roboto font family via Google Fonts.
\end{itemize}

\subsection{Colour Scheme}

\begin{longtable}{p{4.5cm} p{3cm} p{6.5cm}}
\toprule
\textbf{Colour} & \textbf{Hex Code} & \textbf{Usage} \\
\midrule
\endhead
Primary Blue & \texttt{\#1A73E8} & Buttons, links, primary actions \\
Secondary Green & \texttt{\#34A853} & Success states, sensor connected \\
Error Red & \texttt{\#EA4335} & Errors, low glucose alerts \\
Warning Yellow & \texttt{\#FBBC04} & Warnings, activity events \\
Glucose Normal & \texttt{\#66BB6A} & In-range glucose (70--180 mg/dL) \\
Glucose High & \texttt{\#FFA726} & High glucose (180--250 mg/dL) \\
Glucose Very High & \texttt{\#E53935} & Very high glucose ($>$ 250 mg/dL) \\
Glucose Low & \texttt{\#EF5350} & Low glucose ($<$ 70 mg/dL) \\
\bottomrule
\end{longtable}

\subsection{Dark Mode}

A full dark theme is implemented with dedicated background, surface, card, and text colours:
\begin{itemize}
    \item Background: \texttt{\#0F0F0F}
    \item Surface: \texttt{\#1C1C1C}
    \item Cards: \texttt{\#2A2A2A}
    \item Primary text: \texttt{\#FFFFFF}
    \item Secondary text: \texttt{\#E8EAED}
\end{itemize}

%==============================================================================
\section{Glucose Chart Implementation}
%==============================================================================

The \texttt{GlucoseChart} widget uses the \texttt{fl\_chart} library to render an interactive line chart:

\begin{itemize}
    \item \textbf{Y-axis:} Glucose values from 40 to 400~mg/dL with grid lines every 50~mg/dL.
    \item \textbf{X-axis:} Time-based, with labels formatted as \texttt{HH:mm} (hours and minutes).
    \item \textbf{Threshold Lines:} Horizontal dashed lines at 70~mg/dL (low threshold) and 180~mg/dL (high threshold), with a green-shaded target zone between them.
    \item \textbf{Data Points:} Connected by a smooth bezier curve; colour-coded based on glucose status (green for normal, orange for high, red for low/very high).
    \item \textbf{Time Range Selector:} Pill-shaped buttons for 3h, 6h, 12h, and 24h views.
    \item \textbf{Touch Interaction:} Tooltip shows exact value and timestamp on tap.
\end{itemize}

%==============================================================================
\section{Firestore Security Rules}
%==============================================================================

Cloud Firestore security rules enforce role-based access control:

\subsection{Helper Functions}
\begin{itemize}
    \item \texttt{isAuthenticated()} -- Checks \texttt{request.auth != null}.
    \item \texttt{isOwner(userId)} -- Verifies the requesting user owns the document.
    \item \texttt{isDoctor()} -- Reads the user's Firestore document to check \texttt{role == 'doctor'}.
    \item \texttt{isDoctorOfPatient(patientId)} -- Verifies the doctor is assigned to the patient via the \texttt{doctorId} field.
\end{itemize}

\subsection{Collection Rules}

\begin{longtable}{p{3.5cm} p{2cm} p{8.5cm}}
\toprule
\textbf{Collection} & \textbf{Operation} & \textbf{Rule} \\
\midrule
\endhead
\texttt{users} & Read & Owner OR any authenticated doctor \\
\texttt{users} & Create & Owner only (email must match auth token) \\
\texttt{users} & Update & Owner OR doctor (only \texttt{doctorId} field) \\
\texttt{users} & Delete & Owner only \\
\midrule
\texttt{glucoseReadings} & Read & Patient (own data) OR assigned doctor \\
\texttt{glucoseReadings} & Write & Patient only (own data) \\
\midrule
\texttt{diabetesEvents} & Read & Patient (own data) OR assigned doctor \\
\texttt{diabetesEvents} & Write & Patient only (own data) \\
\bottomrule
\end{longtable}

%==============================================================================
\section{Backend API (FastAPI)}
%==============================================================================

\subsection{Overview}

The FastAPI backend serves as a proxy and orchestration layer for:
\begin{itemize}
    \item Dexcom OAuth2 authentication flow.
    \item Background glucose data synchronisation.
    \item AI-powered glucose analysis via LLM integration.
    \item Rate limiting (60,000 requests/hour to Dexcom API).
\end{itemize}

\subsection{Configuration}

The backend uses \texttt{pydantic-settings} to load configuration from environment variables (or a \texttt{.env} file):

\begin{longtable}{p{5cm} p{3cm} p{6cm}}
\toprule
\textbf{Variable} & \textbf{Default} & \textbf{Description} \\
\midrule
\endhead
\texttt{APP\_NAME} & DMS Backend API & Application name \\
\texttt{DEBUG} & False & Enable debug mode \\
\texttt{DEXCOM\_CLIENT\_ID} & (required) & Dexcom OAuth client ID \\
\texttt{DEXCOM\_CLIENT\_SECRET} & (required) & Dexcom OAuth client secret \\
\texttt{DEXCOM\_ENVIRONMENT} & sandbox & \texttt{sandbox}, \texttt{us}, \texttt{eu}, \texttt{jp} \\
\texttt{FIREBASE\_PROJECT\_ID} & (required) & Firebase project ID \\
\texttt{DATABASE\_URL} & (optional) & PostgreSQL connection string \\
\texttt{ENCRYPTION\_KEY} & (required) & AES key for token encryption \\
\bottomrule
\end{longtable}

\subsection{API Endpoints}

\begin{longtable}{p{2cm} p{5.5cm} p{6.5cm}}
\toprule
\textbf{Method} & \textbf{Path} & \textbf{Description} \\
\midrule
\endhead
\texttt{GET} & \texttt{/health} & Service health check \\
\texttt{POST} & \texttt{/dexcom/auth} & Initiate Dexcom OAuth2 flow \\
\texttt{GET} & \texttt{/dexcom/callback} & OAuth2 callback handler \\
\texttt{GET} & \texttt{/dexcom/readings} & Fetch glucose readings \\
\texttt{POST} & \texttt{/sync/trigger} & Trigger manual data sync \\
\texttt{GET} & \texttt{/sync/status} & Check sync job status \\
\bottomrule
\end{longtable}

\subsection{Error Handling}

Custom exception handlers are registered for:
\begin{itemize}
    \item \texttt{DexcomAPIError} -- Returns structured error response with Dexcom-specific details.
    \item \texttt{RateLimitError} -- Returns HTTP~429 with rate limit message.
\end{itemize}

\subsection{Background Jobs}

The \texttt{APScheduler} runs background tasks for:
\begin{itemize}
    \item Periodic glucose data synchronisation from Dexcom to Firestore.
    \item Token refresh for expiring OAuth2 tokens.
\end{itemize}

%==============================================================================
\section{Feature Details}
%==============================================================================

\subsection{Patient Features}

\subsubsection{Dashboard}
The main patient screen displays:
\begin{itemize}
    \item Current glucose value with trend arrow and colour-coded status.
    \item Interactive glucose chart with 3h/6h/12h/24h time range selectors.
    \item Quick statistics: average glucose, time in range, min/max values.
    \item Sensor connection status indicator.
    \item Quick action buttons for adding events and accessing the AI assistant.
\end{itemize}

\subsubsection{Event Logging}
Patients can log the following event types:
\begin{itemize}
    \item \textbf{Insulin} -- Dose in units with type (rapid-acting, long-acting).
    \item \textbf{Meal} -- Carbohydrate count in grams.
    \item \textbf{Activity} -- Duration in minutes with intensity level.
    \item \textbf{Blood Glucose} -- Manual finger-prick reading.
    \item \textbf{Note} -- Free-text observation.
\end{itemize}

\subsubsection{AI Assistant}
An AI-powered chat interface that:
\begin{itemize}
    \item Analyses current and recent glucose patterns.
    \item Provides personalised recommendations and warnings.
    \item Renders responses in Markdown format.
    \item Always includes a medical disclaimer.
    \item Falls back to rule-based analysis when the backend is unavailable.
\end{itemize}

\subsubsection{History}
Historical data view with:
\begin{itemize}
    \item Daily, weekly, and monthly aggregations.
    \item Glucose charts for each time period.
    \item Event timeline overlaid on glucose data.
\end{itemize}

\subsubsection{Health Connections}
Manage doctor--patient relationships:
\begin{itemize}
    \item View assigned doctor's information.
    \item Send connection requests.
    \item Manage data sharing permissions.
\end{itemize}

\subsection{Doctor Features}

\subsubsection{Patient Dashboard}
The doctor's main screen provides:
\begin{itemize}
    \item Searchable patient list fetched from Firestore.
    \item Patient count statistics and filter indicators.
    \item ``Add Patient'' button to connect with patients by email.
    \item Pull-to-refresh for real-time updates.
    \item Alerts tab showing critical patient notifications.
    \item Profile tab with doctor information and settings.
\end{itemize}

\subsubsection{Patient Detail View}
When a doctor selects a patient:
\begin{itemize}
    \item Patient information card (name, email, diabetes type).
    \item Glucose chart with the same time range selectors as the patient view.
    \item Statistics cards (average, in-range percentage, min/max).
    \item Recent events list with icons and relative timestamps.
    \item Data loaded from Firestore (glucose readings + diabetes events).
    \item Pull-to-refresh for updated data.
\end{itemize}

\subsubsection{Add Patient Flow}
\begin{enumerate}
    \item Doctor enters patient's email address in a dialog.
    \item System searches Firestore for a user with that email.
    \item Validates the user exists and has the \texttt{'patient'} role.
    \item Checks for duplicate assignments.
    \item Sets the patient's \texttt{doctorId} field via \texttt{assignPatientToDoctor()}.
    \item Refreshes the patient list automatically.
\end{enumerate}

%==============================================================================
\section{Application Constants}
%==============================================================================

\begin{longtable}{p{5.5cm} p{3cm} p{5.5cm}}
\toprule
\textbf{Constant} & \textbf{Value} & \textbf{Description} \\
\midrule
\endhead
\texttt{glucoseLowThreshold} & 70.0 & Hypoglycaemia boundary (mg/dL) \\
\texttt{glucoseHighThreshold} & 180.0 & Hyperglycaemia boundary (mg/dL) \\
\texttt{glucoseVeryHighThreshold} & 250.0 & Severe hyperglycaemia (mg/dL) \\
\texttt{glucoseTargetMin} & 70.0 & Target range lower bound \\
\texttt{glucoseTargetMax} & 180.0 & Target range upper bound \\
\texttt{chartTimeRanges} & [3, 6, 12, 24] & Available chart durations (hours) \\
\texttt{rolePatient} & \texttt{'patient'} & Patient role identifier \\
\texttt{roleDoctor} & \texttt{'doctor'} & Doctor role identifier \\
\bottomrule
\end{longtable}

%==============================================================================
\section{Dependencies}
%==============================================================================

\subsection{Flutter Dependencies}

\begin{longtable}{p{5.5cm} p{2cm} p{6.5cm}}
\toprule
\textbf{Package} & \textbf{Version} & \textbf{Purpose} \\
\midrule
\endhead
\texttt{go\_router} & 17.0.0 & Declarative URL-based navigation \\
\texttt{provider} & 6.1.2 & Reactive state management \\
\texttt{fl\_chart} & 1.1.1 & Interactive line charts for glucose data \\
\texttt{google\_fonts} & 6.2.1 & Roboto font family (Material Design 3) \\
\texttt{intl} & 0.20.2 & Date/time formatting and internationalisation \\
\texttt{shared\_preferences} & 2.2.3 & Key-value local storage \\
\texttt{firebase\_core} & 4.2.1 & Firebase SDK initialisation \\
\texttt{firebase\_auth} & 6.1.2 & Email/password authentication \\
\texttt{cloud\_firestore} & 6.1.0 & Cloud Firestore database access \\
\texttt{dexcom} & 1.0.7 & Dexcom API wrapper (session init only) \\
\texttt{http} & 1.2.0 & HTTP client for REST API calls \\
\texttt{flutter\_local\_notifications} & 18.0.1 & Local push notification delivery \\
\texttt{timezone} & 0.10.0 & Timezone-aware scheduling \\
\texttt{flutter\_markdown} & 0.7.6 & Markdown rendering for AI responses \\
\bottomrule
\end{longtable}

\subsection{Backend Dependencies}

\begin{longtable}{p{4.5cm} p{2cm} p{7.5cm}}
\toprule
\textbf{Package} & \textbf{Version} & \textbf{Purpose} \\
\midrule
\endhead
\texttt{fastapi} & 0.115.0 & Async Python web framework \\
\texttt{uvicorn} & 0.32.0 & ASGI server with HTTP/2 support \\
\texttt{pydantic} & 2.9.2 & Data validation and serialisation \\
\texttt{pydantic-settings} & 2.6.0 & Environment-based configuration \\
\texttt{httpx} & 0.27.2 & Async HTTP client with HTTP/2 \\
\texttt{firebase-admin} & 6.5.0 & Server-side Firebase SDK \\
\texttt{sqlalchemy} & 2.0.35 & ORM for PostgreSQL \\
\texttt{psycopg2-binary} & 2.9.9 & PostgreSQL driver \\
\texttt{alembic} & 1.13.1 & Database schema migrations \\
\texttt{cryptography} & 44.0.0 & AES encryption for OAuth tokens \\
\texttt{apscheduler} & 3.10.4 & Background job scheduling \\
\bottomrule
\end{longtable}

%==============================================================================
\section{Build and Run Instructions}
%==============================================================================

\subsection{Prerequisites}

\begin{itemize}
    \item Flutter SDK 3.9+ with Dart 3.9+
    \item Android SDK (API 36 recommended) or Xcode for iOS
    \item Firebase project configured with Auth and Firestore
    \item Python 3.12+ for the backend
    \item PostgreSQL (optional, for token storage)
    \item Docker (optional, for containerised backend)
\end{itemize}

\subsection{Flutter Application}

\begin{lstlisting}[style=yamlstyle]
# Install dependencies
cd DMS
flutter pub get

# Run on connected device / emulator
flutter run

# Build release APK
flutter build apk --release

# Build for Windows
flutter build windows --release
\end{lstlisting}

\subsection{Backend Server}

\begin{lstlisting}[style=yamlstyle]
# Set up virtual environment
cd backend
python -m venv venv
source venv/bin/activate  # Linux/macOS
.\venv\Scripts\activate   # Windows

# Install dependencies
pip install -r requirements.txt

# Configure environment
cp .env.example .env
# Edit .env with your credentials

# Run database migrations
alembic upgrade head

# Start the server
uvicorn app.main:app --reload --port 8000
\end{lstlisting}

\subsection{Firebase Setup}

\begin{enumerate}
    \item Create a Firebase project at \url{https://console.firebase.google.com}.
    \item Enable Authentication with Email/Password sign-in method.
    \item Create a Cloud Firestore database.
    \item Run \texttt{flutterfire configure} to generate \texttt{firebase\_options.dart}.
    \item Deploy security rules: \texttt{firebase deploy --only firestore:rules --project <project-id>}.
\end{enumerate}

\subsection{Demo Credentials}

\begin{longtable}{p{4cm} p{5cm} p{5cm}}
\toprule
\textbf{Role} & \textbf{Email} & \textbf{Notes} \\
\midrule
Mock Patient & \texttt{mocked@test.pl} & Auto-generates 7 days of demo data \\
Real Patient & Any email & Requires Dexcom sensor for real data \\
Doctor & Email with ``doctor'' & Sees assigned patients' data \\
\bottomrule
\end{longtable}

%==============================================================================
\section{Known Issues and Limitations}
%==============================================================================

\begin{enumerate}
    \item \textbf{Dexcom Package Bug:} The \texttt{dexcom} Dart package (v1.0.7) fails to parse timestamps with positive UTC offsets (e.g.\ \texttt{+0100}). This is bypassed by making direct HTTP calls to the Dexcom Share API.
    \item \textbf{Unofficial API:} The Dexcom Share API is not officially supported and may change without notice.
    \item \textbf{Firestore Indexes:} Composite queries (e.g.\ timestamp range + patientId) may require creating indexes in the Firebase Console. Currently, time-range filtering is done client-side.
    \item \textbf{Rate Limiting:} The Dexcom API allows 60,000 requests per hour. Excessive polling may hit this limit.
    \item \textbf{Android Emulator Networking:} The Android emulator (API~36) may lose internet connectivity, requiring a cold boot or WiFi toggle to restore.
\end{enumerate}

%==============================================================================
\section{Application Screenshots}
%==============================================================================

% --- Group 1: Authentication Screens ---
\begin{figure}[htbp]
    \centering
    \begin{minipage}{0.3\textwidth}
        \centering
        \includegraphics[width=\linewidth]{login_page.png}
        \caption{Login Screen}
    \end{minipage}
    \hfill
    \begin{minipage}{0.3\textwidth}
        \centering
        \includegraphics[width=\linewidth]{registration.png}
        \caption{Registration}
    \end{minipage}
    \hfill
    \begin{minipage}{0.3\textwidth}
        \centering
        \includegraphics[width=\linewidth]{profile-patient.png}
        \caption{Patient: Profile}
    \end{minipage}
\end{figure}

% --- Group 2: Patient Features ---
\begin{figure}[htbp]
    \centering
    \begin{minipage}{0.3\textwidth}
        \centering
        \includegraphics[width=\linewidth]{main_patient2.png}
        \caption{Patient: Main Screen}
    \end{minipage}
    \hfill
    \begin{minipage}{0.3\textwidth}
        \centering
        \includegraphics[width=\linewidth]{history_patient.png}
        \caption{Patient: History}
    \end{minipage}
    \hfill
    \begin{minipage}{0.3\textwidth}
        \centering
        \includegraphics[width=\linewidth]{alerts_patient.png}
        \caption{Patient: Alerts}
    \end{minipage}
\end{figure}

% --- Group 3: Patient Views and Integrations ---
\begin{figure}[htbp]
    \centering
    \begin{minipage}{0.3\textwidth}
        \centering
        \includegraphics[width=\linewidth]{main_patient.png}
        \caption{Patient: Main Screen}
    \end{minipage}
    \hfill
    \begin{minipage}{0.3\textwidth}
        \centering
        \includegraphics[width=\linewidth]{main_patient3.png}
        \caption{Patient: Glucose Charts View}
    \end{minipage}
    \hfill
    \begin{minipage}{0.3\textwidth}
        \centering
        \includegraphics[width=\linewidth]{health_connections_patient.png}
        \caption{Health Connections}
    \end{minipage}
\end{figure}

\begin{figure}[htbp]
    \centering
    \begin{minipage}{0.3\textwidth}
        \centering
        \includegraphics[width=\linewidth]{ai_assistant.png}
        \caption{Patient: AI Assistant}
    \end{minipage}
    \hfill
    \begin{minipage}{0.3\textwidth}
        \centering
        \includegraphics[width=\linewidth]{demo_analysis.png}
        \caption{Doctor: Profile}
    \end{minipage}
    \hfill
\end{figure}

% --- Group 4: Doctor Views ---
\begin{figure}[htbp]
    \centering
    \begin{minipage}{0.3\textwidth}
        \centering
        \includegraphics[width=\linewidth]{main_doctor.png}
        \caption{Doctor: Main Screen}
    \end{minipage}
    \hfill
    \begin{minipage}{0.3\textwidth}
        \centering
        \includegraphics[width=\linewidth]{profile_doctor.png}
        \caption{Doctor: Profile}
    \end{minipage}
    \hfill
    \begin{minipage}{0.3\textwidth}
        \centering
        \includegraphics[width=\linewidth]{alerts_doctor.png}
        \caption{Doctor: Alerts + Notifications}
    \end{minipage}
\end{figure}

%==============================================================================
\section{Future Roadmap}
%==============================================================================

\begin{enumerate}
    \item \textbf{Push Notifications} -- Firebase Cloud Messaging for real-time critical alerts.
    \item \textbf{Data Export} -- PDF/CSV export of glucose history and statistics.
    \item \textbf{Bluetooth Integration} -- Direct BLE connectivity with CGM sensors.
    \item \textbf{Google/Apple Sign-In} -- Social authentication providers.
    \item \textbf{Unit Switching} -- Real-time mg/dL $\leftrightarrow$ mmol/L conversion.
    \item \textbf{Doctor Notes} -- In-app notes and recommendations from doctor to patient.
    \item \textbf{QR Code Sharing} -- Quick patient--doctor pairing via QR codes.
    \item \textbf{Offline Support} -- Full offline mode with local database and sync-on-reconnect.
    \item \textbf{Unit Tests} -- Comprehensive test coverage for providers, services, and widgets.
    \item \textbf{Internationalisation} -- Multi-language support (Polish, English, German).
\end{enumerate}

% ==============================================================================
\end{document}
% ==============================================================================
